\begin{pdfdiagram}[x=1cm,y=1cm]
  % 取指层。
  \node[hw state, minimum width=1.35cm] (pc) at (0.7,6.2) {PC};
  \node[hw compute, minimum width=2.0cm] (imem) at (3.0,6.2) {指令存储};
  \node[hw state, minimum width=1.35cm] (ir) at (5.4,6.2) {IR};
  \node[hw control, minimum width=1.9cm] (ctrl) at (10.9,6.2) {控制器};
  \draw[hw bus] (pc) -- (imem);
  \draw[hw bus] (imem) -- (ir);
  \draw[hw bus] (ir) -- (ctrl);
  \node[hw label] at (1.85,6.55) {取指地址};
  \node[hw label] at (4.2,6.55) {16-bit 编码};
  \node[hw label] at (8.15,6.55) {opcode};

  % 字段分流先集中，再送往各目的地，避免三条长线叠在 IR 下方。
  \node[hw compute, minimum width=2.4cm] (fields) at (5.1,4.55) {IR 字段分流\\off / imm / r,s,d};
  \draw[hw bus] (ir.south) -- (fields.north);

  % PC 更新层。
  \node[hw compute, minimum width=1.75cm] (add1) at (0.5,4.55) {+1 加法器};
  \node[hw compute, minimum width=1.75cm] (addoff) at (2.65,4.55) {+off 加法器};
  \node[hw compute, minimum width=1.8cm] (pmux) at (1.55,2.8) {pc mux};
  \draw[hw bus] (pc.south) -- (0.7,5.35) -| (add1.north);
  \draw[hw bus] (0.7,5.35) -- (2.65,5.35) -- (addoff.north);
  \draw[hw bus] (fields.west) -- (addoff.east);
  \draw[hw bus] (add1.south) -- (0.5,3.65) -| (pmux.north west);
  \draw[hw bus] (addoff.south) -- (2.65,3.65) -| (pmux.north east);
  \draw[hw bus] (pmux.west) -- (-0.35,2.8) -- (-0.35,6.2) -- (pc.west);
  \node[hw label] at (1.55,3.83) {顺序值 / 分支目标};

  % 译码、执行、访存和写回层。
  \node[hw state, minimum width=1.9cm] (rf) at (4.0,1.05) {寄存器堆};
  \node[hw compute, minimum width=2.5cm] (ext) at (5.1,2.9) {立即数扩展\\imm8 零扩 / off4 符号扩};
  \node[hw compute, minimum width=1.7cm] (muxa) at (7.25,2.15) {mux A};
  \node[hw compute, minimum width=1.7cm] (muxb) at (7.25,0.45) {mux B};
  \node[hw compute, minimum width=2.0cm] (alu) at (9.45,1.3) {ALU / 移位器};
  \node[hw compute, minimum width=1.9cm] (wb) at (11.85,2.15) {写回 mux};
  \node[hw memory, minimum width=2.25cm] (dmem) at (11.85,0.05) {数据存储 / MMIO};
  \draw[hw bus] (fields.south) -- (ext.north);
  \draw[hw bus] (fields.south west) -- (4.0,3.72) -- (rf.north);
  \draw[hw bus] (rf.east) -- (6.1,1.05) |- (muxa.west);
  \draw[hw bus] (6.1,1.05) |- (muxb.west);
  \draw[hw bus] (ext.east) -- (6.45,2.9) -- (muxa.north west);
  \draw[hw bus] (ext.south) -- (5.1,2.35) -- (6.55,2.35) -- (6.55,0.8) -- (muxb.north west);
  \draw[hw bus] (muxa.east) -- (alu.north west);
  \draw[hw bus] (muxb.east) -- (alu.south west);
  \draw[hw bus] (alu.north east) -- (wb.west);
  \node[hw label] at (10.65,2.03) {ALU 结果};
  \draw[hw bus] (alu.east) -- (10.65,1.3) -- (10.65,0.05) -- (dmem.west);
  \node[hw label, anchor=east] at (10.48,0.72) {地址};
  \draw[hw bus] (dmem.north) -- (11.85,1.18) -- (wb.south);
  \node[hw label, anchor=west] at (12.02,1.05) {读数据};

  % 三条回程各占一个高度，不与前向数据通路共线。
  \draw[hw bus] (rf.south) -- (4.0,-0.2) -- (10.45,-0.2) -- (dmem.west);
  \node[hw label] at (7.2,0.02) {存储写数据};
  \draw[hw bus] (wb.east) -- (12.95,2.15) -- (12.95,-0.85) -- (4.0,-0.85) -- (rf.south);
  \node[hw label] at (8.6,-0.63) {写回数据 + 写地址 d};
  \draw[hw bus] (alu.south) -- (9.45,-1.5) -- (1.55,-1.5) -- (pmux.south);
  \node[hw label] at (5.5,-1.28) {Zero 标志};

  % 控制总线沿最上方走廊扇出。
  \draw[hw return] (ctrl.north) -- (10.9,7.2) -- (6.7,7.2);
  \node[hw label] at (8.8,7.42) {控制总线};
  \draw[hw return] (6.9,7.2) -- (6.9,2.65) -- (muxa.north);
  \draw[hw return] (7.45,7.2) -- (7.45,1.3) -- (alu.west);
  \draw[hw return] (10.4,7.2) -- (10.4,2.65) -- (wb.north west);
  \draw[hw return] (10.9,7.2) -- (13.2,7.2) -- (13.2,0.05) -- (dmem.east);
\end{pdfdiagram}
